\section{Metodología}

\subsection{Diseño de investigación}

La investigación adopta un diseño cuantitativo, comparativo y reproducible por simulación. La unidad independiente es el mundo experimental, identificado por escenario, variante y semilla. El método es un tratamiento repetido sobre ese mundo, no una réplica independiente. Robots, cargas, contactos y muestras temporales quedan anidados dentro de la ejecución. Esta estructura permite comparar métodos sobre condiciones iniciales idénticas sin inflar el tamaño muestral mediante filas correlacionadas.

% ====================================================================
% Cadena metodologica del TFM
% ====================================================================
\begin{figure}[H]
\centering
\resizebox{\textwidth}{!}{%
\begin{tikzpicture}[
  font=\sffamily\small,
  stage/.style={draw=VIUDark!70,fill=VIULight,rounded corners=4pt,
    align=center,text width=3.05cm,minimum height=1.25cm,inner sep=6pt},
  key/.style={draw=VIUOrange,fill=VIUOrange!10,rounded corners=4pt,
    align=center,text width=3.05cm,minimum height=1.25cm,inner sep=6pt,
    font=\sffamily\small\bfseries},
  audit/.style={draw=VIUDark!55,fill=VIUWhite,rounded corners=4pt,
    align=center,text width=3.05cm,minimum height=0.9cm,inner sep=5pt,
    font=\sffamily\scriptsize},
  flow/.style={-{Latex[length=2.4mm]},line width=0.9pt,draw=VIUOrange!90!black},
  soft/.style={-{Latex[length=2.0mm]},line width=0.7pt,draw=VIUDark!55}
]
\node[stage] (q) at (0,0) {\textbf{Pregunta}\\coaliciones factibles};
\node[stage] (m) at (3.6,0) {\textbf{Modelo}\\robots\\cargas\\grafo};
\node[stage] (t) at (7.2,0) {\textbf{Teor\'ia}\\utilidad, Smith, cotas};
\node[key]   (e) at (10.8,0) {\textbf{Experimentos}\\SP1--SP8 pareados};
\node[stage] (r) at (14.4,0) {\textbf{Lectura}\\regímenes y límites};

\node[audit] (a1) at (0,-1.95) {fuente bibliografica\\o problema industrial};
\node[audit] (a2) at (3.6,-1.95) {parametros, semillas\\y escenarios};
\node[audit] (a3) at (7.2,-1.95) {resultado algebraico\\o hipótesis};
\node[audit] (a4) at (10.8,-1.95) {CSV, auditorías\\y contraste Holm};
\node[audit] (a5) at (14.4,-1.95) {afirmaci\'on acotada\\sin sobreafirmar};

\foreach \from/\to in {q/m,m/t,t/e,e/r}
  \draw[flow] (\from) -- (\to);
\foreach \from/\to in {q/a1,m/a2,t/a3,e/a4,r/a5}
  \draw[soft] (\from.south) -- (\to.north);

\node[draw=VIUOrange!75,fill=VIUOrange!6,rounded corners=4pt,
  align=center,text width=12.8cm,inner sep=6pt,font=\sffamily\scriptsize,
  below=0.62cm of a3] (rule)
  {Regla de cierre: una afirmación fuerte entra en la memoria solo si apunta a una prueba, una cita, un CSV canónico o una limitación explícita.};
\draw[soft] (a4.south) -- (rule.north);
\end{tikzpicture}}
\caption[Cadena metodológica del TFM]{Cadena metodológica del TFM. La investigación conecta pregunta, modelo, teoría, tratamientos y métricas; la validación no se declara como prueba formal cuando depende de efectos discretos o físicos no cerrados analíticamente.}
\viuownsource
\label{fig:metodologia_investigacion}
\end{figure}


El procedimiento general comprende cinco operaciones:
\begin{enumerate}[leftmargin=2em]
  \item generar y fijar mundos con robots, cargas, grafo, obstáculos, batería, demanda física y eventos;
  \item separar semillas de desarrollo, ajuste, validación, entrenamiento y prueba;
  \item ejecutar cada tratamiento sobre el mismo mundo y preservar éxitos, fallos, excepciones y timeouts;
  \item calcular métricas físicas, operativas y computacionales con semántica y unidades explícitas;
  \item decidir cada afirmación mediante efecto, intervalo de confianza, multiplicidad y límite de interpretación.
\end{enumerate}

\subsection{Unidad experimental, emparejamiento y dependencia}

Para una comparación entre métodos $a$ y $b$, la observación básica es la diferencia dentro del mismo mundo,
\[
  d_w = y_{w,a}-y_{w,b},
\]
donde $w$ identifica escenario, variante y semilla. Las inferencias agregan $d_w$ o modelan el resultado con un efecto de mundo. No se tratan como independientes las mediciones por robot, carga, frame, método o checkpoint cuando proceden del mismo mundo. Las curvas temporales son trayectorias descriptivas; su uso inferencial requiere un modelo longitudinal o un resumen predefinido por ejecución.

\begin{table}[H]
\centering
\footnotesize
\begin{tabularx}{\textwidth}{@{}p{0.20\textwidth}p{0.31\textwidth}X@{}}
\toprule
Elemento & Papel & Regla de análisis \\
\midrule
Mundo & Unidad independiente & Define el cluster y el bloque de emparejamiento. \\
Método o variante & Tratamiento repetido & Se compara dentro del mismo mundo. \\
Robot, carga o contacto & Unidad anidada & No incrementa por sí sola el número de réplicas. \\
Frame o iteración & Medida longitudinal & Se resume por ejecución o se modela con dependencia temporal. \\
Semilla de entrenamiento & Réplica del proceso de aprendizaje & Se conserva como unidad; no se elige la mejor semilla. \\
Timeout o error & Resultado observado & Se preserva; nunca se elimina ni se convierte en dato ausente. \\
\bottomrule
\end{tabularx}
\caption{Jerarquía de unidades y tratamiento de la dependencia.}
\label{tab:variables-diseno-experimental}
\end{table}

\subsection{Comparadores, ajuste y prueba}

La comparación incluye reglas voraces, asignación húngara, subastas o consenso, métodos basados en modelo, referencias centralizadas, variantes propuestas y políticas aprendidas. Las etiquetas \emph{proxy}, \emph{adaptación} o \emph{inspirado en} identifican implementaciones que no reproducen íntegramente la formulación y el alcance de información de la fuente. Un método solo se presenta como reproducción fiel cuando coinciden algoritmo, observación, restricciones y protocolo.

Cuando intervienen aprendizaje o ajuste, el protocolo exige registros de semillas disjuntos para entrenamiento, ajuste, validación y prueba. Los hiperparámetros y campeones se seleccionan con validación. En las campañas confirmatorias con \emph{freeze}, las semillas confirmatorias permanecen selladas hasta la congelación del protocolo. Las políticas finales se evalúan mediante el checkpoint fijado por presupuesto; una semilla fallida se conserva y no se reemplaza. Las referencias con información global se reportan como cotas internas o comparadores, no como soluciones desplegables.

La Tabla~\ref{tab:protocolo-marl} evita presentar un nombre de algoritmo como evidencia suficiente. El comparador aprendido conservado es reproducible en arquitectura y checkpoint, pero no supera sus gates de comportamiento. Por ello participa en la taxonomía de comparación y en la discusión de O5, no en una afirmación favorable de rendimiento.

\begin{table}[H]
\centering
\scriptsize
\begin{tabularx}{\textwidth}{@{}p{0.12\textwidth}p{0.20\textwidth}p{0.18\textwidth}p{0.22\textwidth}X@{}}
\toprule
Comparador & Actor/observación & Crítico y entrenamiento & Réplicas y checkpoints & Resultado de auditoría \\
\midrule
MAPPO--GNN cfg03 & Actor GNN compartido, 3 capas y 64 unidades; observación del vecindario del grafo; ejecución local. & Crítico global solo durante entrenamiento CTDE; PPO recortado; horizonte 4; $\gamma=0{,}95$; CPU. & Semillas 15001--15003; 200.000 transiciones conjuntas por semilla. SHA-256 finales: \texttt{de6de331\ldots}, \texttt{58655bb1\ldots}, \texttt{ac07f856\ldots}. & Checkpoints cargables, pero éxito RAW $=0$ para actor aleatorio, finales y perturbado; éxito tras reparación $=1$. Gate de sensibilidad FAIL: no se autoriza superioridad aprendida. \\
Controles de atribución & Actor aleatorio y logits uniformes sobre los mismos mundos. & Sin entrenamiento; misma ruta de evaluación y reparación. & Controles fijados por la auditoría SP0. & Reproducen RAW $=0$ y reparación $=1$; confirman dominancia del cierre determinista. \\
\bottomrule
\end{tabularx}
\caption{Protocolo del comparador MARL y decisión de promoción. Los hashes completos y los gates se conservan en \texttt{training/champion.yaml} y \texttt{SP0\_AUDIT\_REPORT.md}.}
\label{tab:protocolo-marl}
\end{table}

\subsection{Presupuesto de centralización}

La palabra \emph{distribuida} describe únicamente las etapas cuya actualización y observación satisfacen la vista local declarada. La implementación integrada es un instrumento diagnóstico centralizado que evalúa componentes distribuibles y cierres globales sobre la misma planta; no es un controlador local de extremo a extremo. La Tabla~\ref{tab:presupuesto-centralizacion} hace visible esa deuda de información.

\begin{table}[H]
\centering
\scriptsize
\begin{tabularx}{\textwidth}{@{}p{0.11\textwidth}p{0.25\textwidth}p{0.31\textwidth}X@{}}
\toprule
Etapa & Señal disponible por robot & Operación global del artefacto & Alcance defendible \\
\midrule
A0 & Salud, capacidad y distancia para calcular utilidad individual. & Cuantil global $0{,}55$ y selección sobre toda la flota; registra $N$ mensajes. & Preferencia por robot; no cierre local. \\
A1 & Preferencia y distancia. & Ranking de todos los candidatos, cuórum y slots. & Cierre entero auditado con información global. \\
A2 & Capacidad efectiva y salud. & Ordenación de ociosos y suma global hasta cubrir demanda. & Comprobación global de capacidad; no mercado local cerrado. \\
A3 & Demanda y contribución wrench de contactos. & Vecindario global de altas/bajas de uno o dos robots y asignación de slots. & Búsqueda acotada; coste de candidatos $O(N^2)$ en el peor caso. \\
A4 & Orden HOCBF y wrench realizable. & Unión global de las coaliciones A2--A3 antes de ejecutar. & Reserva dinámica en la planta común; no optimización global de trayectoria. \\
FULL & Mensajes retardados, pérdida y memoria expirable por robot. & Selección del reemplazo entre todos los robots ociosos usando capacidad, salud y posición. & Proxy de recuperación bajo red degradada; no reemplazo vecino extremo a extremo. \\
SP1--SP3 & Payoffs, ocupaciones o precios según la variante. & QR, \emph{closure} y guardia wrench consultan preferencias o candidatos globales. & Aíslan juego continuo y cierre discreto; CLOSED no acredita localidad total. \\
SP4--SP6 & Utilidad de primitivas, barreras y estados físicos por variante. & Prioridades, cierre de contactos o reemplazo pueden ser globales. & Evidencia de movimiento, carga y fallo dentro de cada planta simulada. \\
SP7--SP8 & Descriptores de red o agregados por zona. & La red no cierra el controlador histórico y el evaluador mesoscópico conoce el método. & Evidencia descriptiva o exploratoria, no escalado distribuido observado. \\
\bottomrule
\end{tabularx}
\caption{Presupuesto de centralización. Una fórmula local no convierte en local el umbral, el cierre o la guardia que consume estado global.}
\label{tab:presupuesto-centralizacion}
\end{table}

\begin{landscape}
\scriptsize
\begin{longtable}{p{1.3cm}p{2.6cm}p{3.0cm}p{2.5cm}p{2.5cm}p{2.4cm}p{2.7cm}}
\caption{Generador, mecanismo, cierre, salvaguarda, recuperación y localidad por experimento.}\label{tab:mecanismo-por-experimento}\\
\toprule
Campaña & Generador/planta & Jugadores y acciones & Cierre & Guardia/recuperación & Localidad efectiva & Afirmación autorizada \\
\midrule
\endfirsthead
\toprule Campaña & Generador/planta & Jugadores y acciones & Cierre & Guardia/recuperación & Localidad efectiva & Afirmación autorizada \\
\midrule
\endhead
A0--FULL & Carga Euler--Lagrange planar, cuatro familias. & A0: utilidad y umbral, no juego; capas posteriores: capacidad, contactos y mensajes. & Ranking, cuórum, búsqueda de uno/dos robots y unión A2--A3. & HOCBF; FULL sustituye como máximo un miembro por baja. & Umbral, cierres y selector globales. & Efecto acumulativo del paquete añadido, no causalidad de un componente aislado. \\
SP1 & Asignación homogénea sintética. & Robots reparten preferencia; Smith, replicator y BNN. & QR entero. & Sin recuperación física. & Actualización local por variante; QR canónico global. & Separa preferencia y cierre; no superioridad general de Smith. \\
SP2 & Capacidades heterogéneas sintéticas. & Robots eligen carga con fitness de capacidad. & Closure por capacidad y distancia. & Control uniforme como atribución rival. & Cierre global. & La capacidad importa; el decodificador puede dominar la dinámica. \\
SP3 & Contactos y demanda planar cuasiestática. & Robots eligen carga--slot; primal--dual/poblacionales. & Cierre entero y asignación de slots. & Guardia global añade individuos o pares. & Precios pueden ser locales; guardia consulta todos los libres. & KKT/v-GNE de la relajación y necesidad del residual; CLOSED no se atribuye al juego. \\
SP4 & Uniciclo dinámico reducido. & Robots eligen primitivas espacio--tiempo. & Reserva greedy global de recursos. & Proyección CBF; sin recuperación de coalición. & Exacta o consenso según variante; cierre global. & Resultados de docking en 18 bloques; no transporte ni seguridad funcional. \\
SP5 & Carga Euler--Lagrange desde contactos fijos. & Controladores generan wrench nominal. & Reparto a fuerzas acotadas. & CBF local/preview; RAW--SAFE--EXEC. & Depende de variante; contactos predefinidos. & Transporte simulado; cero colisiones no implica invariancia/hardware. \\
SP6 & Fallos sintéticos durante transporte. & Métodos reparan o conservan coalición. & Reasignación por variante. & Guardia de factibilidad posterior al evento. & Histórica; autoridad global por variante. & Evidencia descriptiva de recuperación simulada. \\
SP7 & Red/sensado sintéticos sobre mundos SP5. & Métodos reciben perfiles de comunicación. & No cierra un controlador de red independiente. & Sin gate físico. & Evaluador conoce la red completa. & Descriptores y asociaciones, no tolerancia confirmatoria. \\
SP8 & Generador mesoscópico dependiente del método. & Políticas agregadas por flota/zona. & Reglas de completitud del evaluador. & Mitigación declarada por método. & Agregados globales/zonales. & Tendencias exploratorias; no coste computacional observado. \\
\bottomrule
\end{longtable}
\end{landscape}

\subsection{Métricas y semántica de ejecución}

Las métricas se agrupan en completitud, factibilidad física, seguridad y recursos. La completitud incluye cargas servidas, llegada y recuperación. La factibilidad reúne déficit de capacidad, residual \emph{wrench}, pose y formación. La seguridad registra colisión por ejecución, duración, despeje y activación de correcciones geométricas. Los recursos comprenden energía aproximada, distancia, mensajes, bytes, memoria observada y tiempo en segundos.

SP1 separa la dinámica continua del cierre discreto. Un horizonte agotado puede coexistir con un cierre válido; por ello, \texttt{continuous\_converged}, validez RAW, \texttt{closure\_success} y éxito final se reportan por separado. La preferencia continua, el $\arg\max$ y QR son etapas distintas. El húngaro expandido se registra como óptimo centralizado, nunca como operación del método distribuido.

SP5 v2 aplica la misma disciplina a la mecánica. RAW es el \emph{wrench} nominal. SAFE es la orden después del filtro de barrera. EXEC es el \emph{wrench} que realizan las fuerzas de contacto saturadas y el único que mueve la carga. La campaña no proyecta posiciones después de integrar. Una colisión EXEC no se convierte en éxito porque la orden SAFE fuera factible, y un timeout sin colisión permanece en el denominador. Los contactos se consideran fijos desde el cierre de SP4; por tanto, SP5 evalúa transporte posterior al acoplamiento y no vuelve a acreditar reclutamiento físico.

\subsection{Análisis por tipo de endpoint}

La prueba y el intervalo dependen del endpoint y de la salida efectivamente generada por cada protocolo:
\begin{itemize}[leftmargin=2em]
  \item \textbf{Resultados binarios.} Se informa la diferencia de riesgos pareada con intervalo del $95\%$. Los contrastes congelados usan McNemar exacto, equivalente al test binomial exacto sobre pares discordantes.
  \item \textbf{Regret y magnitudes continuas.} Las campa\~nas can\'onicas usan diferencias por mundo, \emph{bootstrap} percentil pareado e intervalos del $95\%$; cuando estaba preespecificado se a\~nade Wilcoxon pareado. Media, mediana y CVaR se distinguen y conservan sus unidades.
  \item \textbf{Tiempo a soluci\'on.} Los timeouts permanecen como resultados censurados y nunca se convierten en \'exitos ni en datos ausentes. Cuando una campa\~na no genera un modelo de supervivencia, se reportan tasa de timeout y tiempo truncado de forma descriptiva, sin inferir razones de riesgo ni efectos temporales.
  \item \textbf{Mensajes y bytes.} Se resumen sobre los mismos mundos y se contrastan solo cuando existe una salida pareada expl\'icita. No se atribuyen modelos de conteo, mixtos o log-ratios a campa\~nas que no los produjeron.
  \item \textbf{Asociaciones e interacciones.} Los t\'erminos m\'etodo por r\'egimen, m\'etodo por $\log N$, m\'etodo por $\lambda_2$, din\'amica por fitness y din\'amica por cierre se presentan como confirmatorios \'unicamente si aparecen en una tabla de inferencia can\'onica; en otro caso son descripciones reagrupadas por mundo.
\end{itemize}

Cada resultado confirmatorio debe contener estimación, intervalo del $95\%$, tamaño de efecto, número de mundos, eventos, fallos, timeouts, margen y decisión. Un contraste es no estimable si carece de bloques completos, presenta efecto o intervalo no finito, o no identifica la unidad inferencial. En ese caso no se sustituye la ausencia de estimación por $p=1$ ni por una conclusión nula.

\subsection{Relevancia práctica y multiplicidad}

Los contrastes confirmatorios de SP1 se fijaron antes de abrir las semillas 330000--330899. Las cinco comparaciones Smith--referencia usan diferencias por mundo, intervalos pareados del 95\% y corrección de Holm. La decisión combina efecto, intervalo y alcance; la significación aislada no establece superioridad práctica.

\subsection{Protocolo integrado de comprobaciones físicas}

El protocolo integrado v1.1 de tamaño fijo se preespecificó para resolver una limitación de la escalera modular: SP1--SP8 no comparten una única planta. Cuatro familias de mundo aíslan rotación nominal, escasez, complementariedad de torque y obstáculo con pérdida de paquetes, retardo y \emph{dropout}. Cada mundo se ejecuta con seis prefijos acumulativos A0--FULL y el mismo \texttt{world\_hash}; por ello, el contraste relevante es el cambio dentro del mundo al añadir una comprobación.

El contrato correctivo fija, antes de abrir semillas, 100 mundos independientes por familia y seis tratamientos repetidos por mundo: 400 mundos y 2.400 ejecuciones. No existe checkpoint intermedio ni regla de extensión. Para cada uno de los veinte cambios consecutivos se estima la diferencia de riesgos \emph{después menos antes}; se calcula un IC95\% mediante 4.000 remuestreos pareados con semilla de análisis congelada y se aplica McNemar exacto bilateral. Holm corrige conjuntamente los veinte valores $p$. Un efecto positivo o negativo solo se respalda si el intervalo excluye cero en esa dirección y $p_{\rm Holm}<0{,}05$; cualquier otro resultado es inconcluso o nulo.

Los estados \texttt{accepted\_by\_stage}, \texttt{final\_physical\_success} y \texttt{physical\_false\_positive} se conservan por separado. La primera variable solo acredita lo que el escalón sabe comprobar. La segunda exige todas las comprobaciones estáticas y el resultado dinámico. La tercera identifica la aceptación que se vuelve inválida al ejecutar la carga. La semántica impide que un QP resuelto, un waypoint alcanzado o una reparación local oculten una colisión, un déficit o un fallo de llegada.

La ejecución fue exclusivamente CPU. El manifiesto registra Python~3.13.9, NumPy~2.3.5, 22 procesadores lógicos y cuatro \emph{workers} con un hilo numérico por \emph{worker}. El \emph{dry-run} usó semillas ajenas al conjunto confirmatorio. El registro de semillas se abrió después del manifiesto congelado; se preservaron fallos, hashes y las 2.400 filas únicas. La campaña adaptativa v1 queda supersedida y no se mezcla con estas estimaciones.
\subsection{Validez, reproducibilidad y evidencia visual}

Una campaña se promueve a evidencia canónica cuando dispone de configuración congelada, semillas, hashes, mundos compartidos, tablas de ejecución, esquema válido y manifiesto de cierre. Los fallos y timeouts forman parte del denominador. Las campañas históricas que no cumplen el contrato actual se conservan como evidencia exploratoria o de desarrollo y no se renombran como confirmatorias.

Las secuencias visuales permiten inspeccionar formación, movimiento, contacto y recuperación. No sustituyen la inferencia ni acreditan comportamiento físico con hardware. CoppeliaSim se usa, cuando corresponde, como inspección cualitativa separada; SP1 homogéneo emplea exclusivamente el simulador numérico Python. Las limitaciones incluyen contacto planar simplificado, modelos sintéticos de comunicación, proyección geométrica en algunas campañas y agregación mesoscópica en las flotas mayores.

\subsection{Identidad canónica de la evidencia}

Para evitar que una versión exploratoria sustituya a la evidencia confirmatoria, la memoria congela una identidad por experimento. SP4 v3 (\texttt{SP4\_DOCKING\_GAME\_CONFIRMATORY\_v3}) es la campaña canónica: 18 bloques independientes definidos por semilla y tamaño de flota, observados en seis escenarios, para 108 instancias pareadas y 1.188 ejecuciones de método. Las tasas sobre 108 instancias son descriptivas; la sensibilidad inferencial agrega primero cada endpoint dentro del bloque y aplica un test de signos exacto con Holm. SP4 v4 y sus variantes no promocionadas quedan como evidencia complementaria y no modifican las conclusiones de v3. La escena CoppeliaSim se conserva como \emph{replay} cinemático de trayectorias precalculadas; no ejecuta dinámica independiente, contacto cerrado ni hardware.

La misma regla se aplica a SP8: sus tendencias de calidad son exploratorias y sus límites de tiempo y memoria son declarados o analíticos, no mediciones externas de \emph{wall-clock} o RSS. El protocolo conserva la correspondencia entre resultados, figuras y afirmaciones.

\begin{table}[H]
\centering
\scriptsize
\begin{tabularx}{\textwidth}{@{}p{0.25\textwidth}p{0.26\textwidth}X@{}}
\toprule
Artefacto & Estado & Qué permite afirmar \\
\midrule
SP4 v3 confirmatorio & Canónico & Tasas y contrastes declarados para el juego de docking en el simulador reducido. \\
SP4 v4 y escenas auxiliares & Complementario/no canónico & Diagnóstico adicional; no sustituye v3 ni se mezcla con sus denominadores. \\
CoppeliaSim paired scene & Replay cinemático & Inspección geométrica y visual de trayectorias precalculadas; no validación dinámica independiente ni hardware. \\
HIL & No realizado & No existen entradas/salidas temporizadas, hashes ni datos de sensores/actuadores; no aporta evidencia. \\
AMR físico & No realizado & No existen ensayos, registro de riesgos ni validación de seguridad funcional. \\
SP8 fleet ladder & Exploratorio & Tendencias del generador mesoscópico; sin medición externa de recursos. \\
\bottomrule
\end{tabularx}
\caption{Gobierno de la evidencia y alcance de los artefactos no canónicos.}
\label{tab:identidad-evidencia}
\end{table}

\begin{table}[H]
\centering
\scriptsize
\begin{tabularx}{\textwidth}{@{}p{0.13\textwidth}p{0.17\textwidth}p{0.16\textwidth}p{0.18\textwidth}X@{}}
\toprule
Modalidad & Motor/dinámica & Control y sensores & Entradas/salidas trazables & Alcance \\
\midrule
A0--FULL v1.1 & Python; carga Euler--Lagrange planar. & Selector, cierres, HOCBF y recuperación simulados; estado sintético completo. & YAML, registro de semillas, CSV/Parquet, hashes y manifiesto cerrado. & Evidencia integrada canónica dentro del modelo. \\
SP4 v3 & Python; uniciclo dinámico reducido. & Primitivas, juego y filtros con observación definida por variante. & Configuración, tabla de ejecuciones, trazas, auditoría y sensibilidad por bloques. & Evidencia modular de \emph{docking}; no contacto físico real. \\
CoppeliaSim SP4 v4 & Replay cinemático. & Reproduce trayectorias precalculadas; sin controlador ni sensores independientes. & Escena y trayectorias; no cadena cerrada de entrada--salida física. & Inspección visual, no validación dinámica. \\
HIL & No implementado. & No disponible. & No hay artefacto ni hash de campaña. & Sin evidencia. \\
Hardware AMR & No implementado. & No disponible. & No hay ensayo ni registro de seguridad. & Sin evidencia. \\
\bottomrule
\end{tabularx}
\caption{Modalidades de validación y frontera entre simulación, replay, HIL y hardware. Cero colisiones en una campaña numérica no equivale a seguridad funcional.}
\label{tab:modalidades-validacion}
\end{table}
\subsection{Jerarquía de la evidencia}

A0--FULL decide la hipótesis integradora sobre una sola planta. Las campañas SP1--SP3 explican el paso desde preferencia continua hasta capacidad y contactos; SP4--SP6 aíslan movimiento, carga extendida y recuperación; SP7--SP8 delimitan preguntas sobre red y escala que permanecen descriptivas o exploratorias. Esta jerarquía impide usar una campaña histórica para sostener una afirmación que exige ejecución integrada.

El registro único de la Tabla~\ref{tab:claim-artifact-registry} enlaza cada afirmación con su artefacto y su límite. Su versión completa, incluidos configuración, comando, semillas, manifiesto y SHA-256, se genera en \texttt{docs/generated/claim\_artifact\_registry.csv}; el registro de estimandos y contrastes se conserva en \texttt{contrast\_estimand\_registry.csv}.

\begin{landscape}
\scriptsize
\begin{longtable}{p{1.2cm}p{1.5cm}p{3.5cm}p{7.8cm}p{6.7cm}}
\caption{Registro único afirmación--artefacto y límite de inferencia.}\label{tab:claim-artifact-registry}\\
\toprule
ID & Campaña & Estado & Afirmación autorizada & Límite decisivo \\
\midrule
\endfirsthead
\toprule ID & Campaña & Estado & Afirmación autorizada & Límite decisivo \\
\midrule
\endhead
C-INT & A0--FULL & confirmatoria en simulación planar & La escalera separa cardinalidad, capacidad, wrench, mecánica y recuperación; los efectos dependen de la familia. & Planta planar numérica; selector FULL con candidatos globales; sin fricción de contacto ni hardware. \\
C-SP0 & SP0 & histórica bloqueada para MARL & Existen tres checkpoints MAPPO-GNN trazables de 200.000 transiciones; el éxito cerrado observado procede de la reparación. & RAW no supera controles; separación reparación/QR y varios gates estadísticos fallan. No autoriza superioridad aprendida. \\
C-SP1 & SP1 & confirmatoria modular & Smith-QR mejora a greedy en regret, pero no a todos los controles; el cierre explica parte sustancial del desempeño. & El cierre QR usa preferencias globales; no acredita distribución extremo a extremo. \\
C-SP2 & SP2 & confirmatoria modular & La capacidad efectiva mejora la señal, pero ninguna dinámica domina y uniforme+closure es competitivo. & Capacidad escalar; no certifica fuerza, torque ni contacto. \\
C-SP3 & SP3 & confirmatoria modular & La relajación satisface la referencia QP/KKT y la guardia elimina falsos positivos wrench en el generador. & La guardia es global y domina CLOSED; equivalencia continua no implica optimalidad entera. \\
C-SP4 & SP4 & confirmatoria modular con sensibilidad & Los cinco contrastes conservan dirección y soporte al agregar por 18 bloques independientes. & 18 bloques semilla-flota; 108 instancias no son réplicas independientes; uniciclo dinámico reducido. \\
C-SP5 & SP5 & confirmatoria modular & El filtro CBF reduce colisiones de la variante Hamiltoniana RAW, sin demostrar dominancia global. & Simulación CPU desde contactos fijos; cero colisiones observadas no equivale a seguridad funcional. \\
C-SP6 & SP6 & descriptiva histórica & La campaña mide reemplazo, recuperación y pérdida de carga bajo fallos simulados. & Modelo de fallo y recuperación simulado; no HIL ni robot real. \\
C-SP7 & SP7 & descriptiva & La canalización cuantifica conectividad, pérdida, retardo, sensing y éxito dentro de su generador. & Sin freeze/manifiesto final promovible ni modelo RF; no se eleva a evidencia confirmatoria. \\
C-SP8 & SP8 & exploratoria & El generador mesoscópico describe tendencias de calidad entre 5 y 50.000 AMR. & Tiempo y memoria no fueron medidos con watchdog/RSS común; no demuestra intractabilidad industrial. \\
\bottomrule
\end{longtable}
\end{landscape}

